\subsection{Exercices}

\textit{Pour tous les exercices de calcul, vous pouvez utiliser les valeurs (arrondies) suivantes pour la fonction de répartition de la loi normale standard $\Phi(z) = P(Z \le z)$ :}
\begin{itemize}
    \item $\Phi(0) = 0.5$
    \item $\Phi(1.0) \approx 0.8413$
    \item $\Phi(1.5) \approx 0.9332$
    \item $\Phi(1.96) \approx 0.975$
    \item $\Phi(2.0) \approx 0.9772$
    \item $\Phi(2.5) \approx 0.9938$
    \item $\Phi(3.0) \approx 0.9987$
\end{itemize}
\textit{Et rappelez-vous la propriété de symétrie : $\Phi(-z) = 1 - \Phi(z)$.}

% --- Section 1 : Paramètres de S_n et X_n barre ---

\begin{exercicebox}[Exercice 1 : Paramètres de $\bar{X}_n$]
Soit $X_i$ une suite de v.a. i.i.d. avec $\mu=50$ et $\sigma^2=100$. Soit $n=25$.
Calculez $E[\bar{X}_{25}]$ et $\text{Var}(\bar{X}_{25})$.
\end{exercicebox}

\begin{exercicebox}[Exercice 2 : Paramètres de $S_n$]
Soit $X_i$ une suite de v.a. i.i.d. avec $\mu=10$ et $\sigma=3$. Soit $n=16$.
Calculez $E[S_{16}]$ et l'écart-type $\sigma_{S_{16}}$.
\end{exercicebox}

\begin{exercicebox}[Exercice 3 : Paramètres de $\bar{X}_n$ (bis)]
Soit $X_i$ une suite de v.a. i.i.d. avec $\mu=70$ et $\sigma=10$. Soit $n=400$.
Calculez $E[\bar{X}_{400}]$ et l'écart-type $\sigma_{\bar{X}_{400}}$.
\end{exercicebox}

\begin{exercicebox}[Exercice 4 : Paramètres de $S_n$ (bis)]
Soit $X_i$ une suite de v.a. i.i.d. avec $\mu=0.5$ et $\sigma^2=0.01$. Soit $n=64$.
Calculez $E[S_{64}]$ et $\text{Var}(S_{64})$.
\end{exercicebox}

\begin{exercicebox}[Exercice 5 : Retrouver $\sigma^2$ (via $\bar{X}_n$)]
La moyenne d'échantillon $\bar{X}_n$ de $n=49$ observations a une variance $\text{Var}(\bar{X}_{49}) = 2$.
Quelle est la variance $\sigma^2$ de la population d'origine ?
\end{exercicebox}

\begin{exercicebox}[Exercice 6 : Retrouver $n$ (via $S_n$)]
La somme $S_n$ d'observations i.i.d. a une variance $\text{Var}(S_n) = 300$. La variance de la population est $\sigma^2 = 12$.
Quelle est la taille de l'échantillon $n$ ?
\end{exercicebox}

% --- Section 2 : Calculs de probabilité (Forme $\bar{X}_n$) ---

\begin{exercicebox}[Exercice 7 : CLT pour $\bar{X}_n$ (Queue Droite)]
Soit $X_i$ i.i.d. avec $\mu=100$ et $\sigma=15$. Soit $n=36$.
Calculez $P(\bar{X}_{36} > 105)$.
\end{exercicebox}

\begin{exercicebox}[Exercice 8 : CLT pour $\bar{X}_n$ (Queue Gauche)]
Soit $X_i$ i.i.d. avec $\mu=50$ et $\sigma=8$. Soit $n=64$.
Calculez $P(\bar{X}_{64} < 48)$.
\end{exercicebox}

\begin{exercicebox}[Exercice 9 : CLT pour $\bar{X}_n$ (Intervalle)]
Soit $X_i$ i.i.d. avec $\mu=20$ et $\sigma=5$. Soit $n=100$.
Calculez $P(19 \le \bar{X}_{100} \le 21.25)$.
\end{exercicebox}

\begin{exercicebox}[Exercice 10 : Application $\bar{X}_n$ (Bouteilles)]
Une machine remplit des bouteilles avec $\mu=500$ ml et $\sigma=6$ ml. On prend $n=36$ bouteilles.
Calculez $P(\bar{X}_{36} > 501.5)$.
\end{exercicebox}

\begin{exercicebox}[Exercice 11 : Application $\bar{X}_n$ (Tailles)]
La taille des individus a $\mu=175$ cm et $\sigma=8$ cm. On prend $n=64$ individus.
Calculez $P(\bar{X}_{64} < 173)$.
\end{exercicebox}

\begin{exercicebox}[Exercice 12 : Application $\bar{X}_n$ (Notes)]
Les notes à un examen ont $\mu=70$ et $\sigma=12$. Une classe de $n=36$ étudiants est un échantillon.
Calculez $P(\bar{X}_{36} < 67)$.
\end{exercicebox}

% --- Section 3 : Calculs de probabilité (Forme $S_n$) ---

\begin{exercicebox}[Exercice 13 : CLT pour $S_n$ (Queue Droite)]
Soit $X_i$ i.i.d. avec $\mu=10$ et $\sigma=2$. Soit $n=100$.
Calculez $P(S_{100} > 1020)$.
\end{exercicebox}

\begin{exercicebox}[Exercice 14 : CLT pour $S_n$ (Queue Gauche)]
Soit $X_i$ i.i.d. avec $\mu=5$ et $\sigma=4$. Soit $n=64$.
Calculez $P(S_{64} \le 304)$.
\end{exercicebox}

\begin{exercicebox}[Exercice 15 : CLT pour $S_n$ (Intervalle)]
Soit $X_i$ i.i.d. avec $\mu=2$ et $\sigma=3$. Soit $n=36$.
Calculez $P(S_{36} > 90)$.
\end{exercicebox}

\begin{exercicebox}[Exercice 16 : Application $S_n$ (Ascenseur)]
Un ascenseur transporte $n=49$ personnes. Poids : $\mu=70$kg, $\sigma=14$kg.
Calculez $P(S_{49} > 3528)$.
\end{exercicebox}

\begin{exercicebox}[Exercice 17 : Application $S_n$ (Rendement)]
Le rendement quotidien $X_i$ a $\mu=0.001$ et $\sigma=0.01$. Soit $n=100$.
Calculez $P(S_{100} > 0.2)$.
\end{exercicebox}

\begin{exercicebox}[Exercice 18 : Application $S_n$ (Rendement)]
En utilisant les données de l'exercice 17, calculez $P(S_{100} < 0)$.
\end{exercicebox}

% --- Section 4 : Problèmes Inverses (CLT) ---

\begin{exercicebox}[Exercice 19 : Inverse (Trouver $c$ pour $\bar{X}_n$)]
Soit $X_i$ i.i.d. avec $\mu=50$ et $\sigma=10$. Soit $n=100$.
Trouvez la valeur $c$ telle que $P(\bar{X}_{100} \le c) \approx 0.8413$.
\end{exercicebox}

\begin{exercicebox}[Exercice 20 : Inverse (Trouver $c$ pour $S_n$)]
Soit $X_i$ i.i.d. avec $\mu=10$ et $\sigma=3$. Soit $n=36$.
Trouvez la valeur $c$ telle que $P(S_{36} \le c) \approx 0.0013$.
\end{exercicebox}

\begin{exercicebox}[Exercice 21 : Inverse (Taille d'échantillon $n$)]
Une population a $\mu=0$ et $\sigma=10$.
Quelle taille d'échantillon $n$ faut-il pour que $P(|\bar{X}_n| \le 1) \ge 0.95$ ?
(Indice : $P(-1.96 \le Z \le 1.96) = 0.95$).
\end{exercicebox}

\begin{exercicebox}[Exercice 22 : Inverse (Taille d'échantillon $n$)]
Une population a $\mu=100$ et $\sigma=20$.
Quelle taille d'échantillon $n$ faut-il pour que $P(\bar{X}_n \ge 102) \le 0.0228$ ?
\end{exercicebox}

% --- Section 5 : Proportions (Application de Bernoulli) ---

\begin{exercicebox}[Exercice 23 : Paramètres (Proportion)]
On sonde $n=400$ personnes. La vraie proportion $p$ est $0.25$. On modélise $X_i \sim \text{Bern}(p)$.
\begin{enumerate}
    \item Calculez $\mu = E[X_i]$ et $\sigma^2 = \text{Var}(X_i)$.
    \item Calculez $E[\hat{p}]$ et $\text{Var}(\hat{p})$ (où $\hat{p} = \bar{X}_n$).
\end{enumerate}
\end{exercicebox}

\begin{exercicebox}[Exercice 24 : Calcul (Proportion)]
On sonde $n=100$ personnes. La vraie proportion est $p=0.5$.
Calculez la probabilité que la proportion observée $\hat{p}$ soit supérieure à 0.6, $P(\hat{p} > 0.6)$.
\end{exercicebox}

\begin{exercicebox}[Exercice 25 : Marge d'Erreur (Proportion)]
Un sondage sur $n=1000$ personnes donne un résultat $\hat{p} = 0.54$.
Calculez la marge d'erreur à 95\% (c'est-à-dire $1.96 \times SE_{\hat{p}}$).
\end{exercicebox}